\chapter{Nerves and Senses}\label{ch:g5:nerves}

To survive, \omterm{def:g1:plants-animals:animal}{animals} must detect the world and respond. The \omterm{def:g5:nerves:system}{nervous
system} gathers information from \omterm{def:g1:senses:sense}{sense} \omterm{def:g4:body-systems:organs}{organs}, processes it, and sends
signals to \omterm{def:g3:bones:muscle}{muscles} and other effectors. This chapter maps that pathway
for humans.

\section{The nervous system}

\begin{definition}[Nervous system]\label{def:g5:nerves:system}
The \emph{nervous system}\index{nervous system} includes the
\emph{brain}\index{brain}, \emph{spinal cord}\index{spinal cord} and
\emph{nerves}\index{nerve}. It \omterm{def:g1:senses:sense}{senses} conditions, coordinates responses
and enables learning and memory.
\end{definition}

\begin{definition}[Neuron]\label{def:g5:nerves:neuron}
A \emph{neuron}\index{neuron} (\omterm{def:g5:nerves:system}{nerve} \omterm{def:g5:cell-structure:cell}{cell}) is a \omterm{def:g4:cells:special}{specialised cell} that
carries electrical signals. Many neurons are long so messages can travel
quickly between \omterm{def:g1:my-body:body}{body} parts and the central \omterm{def:g5:nerves:system}{nervous system}.
\end{definition}

\begin{omfigure}
\begin{tikzpicture}[font=\small,
  box/.style={draw, rounded corners, minimum width=2.3cm, minimum height=0.75cm,
    align=center, fill=omDef!12}]
  \node[box] (s) at (0,0) {sense organ};
  \node[box, fill=omThm!12] (b) at (3.5,0) {brain /\\spinal cord};
  \node[box, fill=omMeth!12] (m) at (7.0,0) {muscle or\\gland};
  \draw[-Stealth, thick] (s) -- node[above] {nerve} (b);
  \draw[-Stealth, thick] (b) -- node[above] {nerve} (m);
\end{tikzpicture}

{\small A simplified pathway: detect $\to$ process $\to$ respond.}
\end{omfigure}

\section{Sense organs}

\begin{definition}[Senses]\label{def:g5:nerves:senses}
\omterm{def:g1:senses:sense}{Sense} \omterm{def:g4:body-systems:organs}{organs} detect stimuli:
\begin{itemize}
  \item \emph{Eyes}\index{eye} --- \omterm{def:g4:plants-need:needs}{light} and images.
  \item \emph{Ears}\index{ear} --- sound (and help with balance).
  \item \emph{Skin}\index{skin} --- \omterm{def:g1:senses:five}{touch}, pressure, pain, temperature.
  \item \emph{Tongue}\index{tongue} --- \omterm{def:g1:senses:five}{taste}.
  \item \emph{Nose}\index{nose} --- \omterm{def:g1:senses:five}{smell}.
\end{itemize}
\end{definition}

\begin{example}[Crossing a road]\label{ex:g5:nerves:road}
\omterm{def:g5:nerves:senses}{Eyes} see the traffic \omterm{def:g4:plants-need:needs}{light}; \omterm{def:g5:nerves:senses}{ears} hear engines; the \omterm{def:g5:nerves:system}{brain} decides when
to walk; \omterm{def:g5:nerves:system}{nerves} signal \omterm{def:g1:my-body:parts}{leg} \omterm{def:g3:bones:muscle}{muscles}. Several \omterm{def:g1:senses:sense}{senses} and the \omterm{def:g5:nerves:system}{nervous
system} cooperate.
\end{example}

\begin{proposition}[Stimulus and response]\label{prop:g5:nerves:sr}
A \emph{stimulus}\index{stimulus} is a change that can be detected. A
\emph{response}\index{response} is the action that follows. Fast
responses can protect the \omterm{def:g1:my-body:body}{body} from harm.
\end{proposition}
\admitted

\section{Reflexes}

\begin{definition}[Reflex]\label{def:g5:nerves:reflex}
A \emph{reflex}\index{reflex} is a rapid, automatic response to a
stimulus, often protecting the \omterm{def:g1:my-body:body}{body} (pulling a \omterm{def:g1:my-body:parts}{hand} from heat) without
waiting for conscious thought.
\end{definition}

\begin{example}[Knee jerk and blinking]\label{ex:g5:nerves:examples}
A tap below the kneecap can trigger a kick \omterm{def:g5:nerves:reflex}{reflex} in a medical test.
Blinking when something approaches the \omterm{def:g5:nerves:senses}{eye} protects the surface.
\end{example}

\begin{method}[Tracing a simple response]\label{met:g5:nerves:trace}
\begin{enumerate}
  \item Name the stimulus (for example hot pan).
  \item Name the \omterm{def:g1:senses:sense}{sense} \omterm{def:g4:body-systems:organs}{organ} (\omterm{def:g5:nerves:senses}{skin} sensors).
  \item Note the message path (\omterm{def:g5:nerves:system}{nerves} to \omterm{def:g5:nerves:system}{spinal cord}/\omterm{def:g5:nerves:system}{brain}).
  \item Name the response (\omterm{def:g3:bones:muscle}{muscle} contracts; \omterm{def:g1:my-body:parts}{hand} pulls away).
\end{enumerate}
\end{method}

\begin{omfigure}
\begin{tikzpicture}[font=\footnotesize]
  \draw[thick, omProp, fill=omProp!15] (0,0) circle [radius=0.35];
  \node at (0,-0.7) {hot};
  \draw[-Stealth, thick] (0.5,0.2) -- (2.0,1.0);
  \draw[thick, omDef, fill=omDef!15] (2.5,1.0) circle [radius=0.4];
  \node at (2.5,1.7) {spinal cord};
  \draw[-Stealth, thick] (3.0,0.8) -- (4.5,0.1);
  \draw[thick, omMeth, fill=omMeth!20] (5.0,0) rectangle (6.2,0.5);
  \node at (5.6,-0.5) {muscle moves hand};
\end{tikzpicture}

{\small \omterm{def:g5:nerves:reflex}{Reflex} outline: danger detected, signal processed quickly,
\omterm{def:g3:bones:muscle}{muscles} act.}
\end{omfigure}

\section{Looking after your nervous system}

\begin{example}[Sleep]\label{ex:g5:nerves:sleep}
During sleep the \omterm{def:g5:nerves:system}{brain} consolidates memories and the \omterm{def:g1:my-body:body}{body} recovers.
Tired \omterm{def:g5:nerves:system}{brains} react more slowly --- important for sport and road safety.
\end{example}

\section{Exercises}

\begin{exercise}[$\star$]\label{exo:g5:nerves:1}
Name the three main parts of the \omterm{def:g5:nerves:system}{nervous system} listed in this chapter.
\end{exercise}

\begin{exercise}[$\star$]\label{exo:g5:nerves:2}
What is a \omterm{def:g5:nerves:neuron}{neuron}?
\end{exercise}

\begin{exercise}[$\star$]\label{exo:g5:nerves:3}
Match \omterm{def:g1:senses:sense}{senses} to \omterm{def:g4:body-systems:organs}{organs}: vision, \omterm{def:g1:senses:five}{hearing}, \omterm{def:g1:senses:five}{smell}.
\end{exercise}

\begin{exercise}[$\star$]\label{exo:g5:nerves:4}
What is a stimulus? Give one example.
\end{exercise}

\begin{exercise}[$\star\star$]\label{exo:g5:nerves:5}
What is a \omterm{def:g5:nerves:reflex}{reflex}, and why is it useful?
\end{exercise}

